\documentclass[12pt]{book}
\usepackage{amsmath}
\usepackage{amssymb}
\usepackage{geometry}
\geometry{margin=1in}

\title{Volume 12: Computation, Code, and Simulation Systems \\ Book 03: Data and Derived Artifacts}
\author{James C. Harvey}
\date{December 2025}

\begin{document}
\maketitle
\tableofcontents

\chapter{Data as Derived Geometry}

\section*{Abstract}
Data artifacts in SDT are not raw measurements; they are derived geometric outputs from SDT operators.
This book organizes the derived tables, datasets, and validation artifacts as an explicit layer of the
theory.

\section{Derived Tables}
Computed tables are the numeric projection of curvature, circulation, and occlusion across regimes.

\section{Validation Artifacts}
Derived artifacts encode benchmark alignment and form the formal evidence base for SDT claims.

\chapter{Data Provenance}

\section*{Abstract}
Provenance records the derivation chain that produced each dataset. This chapter defines the provenance
schema used in SDT outputs.

\section{Provenance Fields}
Each dataset records: constants, geometry inputs, operator versions, and timestamped derivation hashes.

\chapter{Formats and Accessibility}

\section*{Abstract}
Data is stored in transparent formats with explicit units. This chapter documents formats and usage.

\section{Format Guidance}
Tables are stored with unit headers and explicit metadata to preserve dimensional clarity.
\end{document}
